\scanpage{36}\mspage{35}
% L1–4
Considérons des $X_i\in M_A$ \unsure{mutuellement} modérés, et les opérations $\unsure{\complement|C}_A,\wedge,\vee$ \deleted{i.e.} \unsure{i.e|(i.e} des $X_i\subset A$ mut. mod (et mod p.r. à $A$ \unsure{autom.}) on regarde les op.
% L5–7
\deleted{$\complement_A\ \wedge,\vee$ \unsure{et} (binaires)\unsure{,} \unsure{$\overline{\phantom X}$}} \deleted{\inserted{binaires \unsure{, $\overline{\phantom X}$}}} $\complement_{M_A}$ (unaire) $\wedge,\vee$ (binaires) et on adjoint $0_{\deleted{\unsure{A}}M_A},1_{\deleted{\unsure{A}}M_A}$, on voit que c'est pareil que
% L8–9
\deleted{\unsure{de l'adjoindre|d'y adjoindre}} de prendre $\complement_M,\wedge,\vee,0_M,1_M\quad\overline{\phantom X}$ et de remultiplier par $A$. \deleted{\unsure{D}}Mais
% L10–12
alors on voit que l'on trouve un idéal d'une alg. de Boole, c'est une alg. de Boole \ldots.

% L13–14
\deleted{\unsure{7|5}} Il faut prouver que
\diagram{\begin{tikzcd}[ampersand replacement=\&,row sep=1.4em,column sep=1em]
\& \underbrace{\{X,Y\}\ \text{mod}}\arrow[d,Rightarrow]\\
\deleted{\unsure{\Longrightarrow}}\ \{\complement_{M_A}X,\complement_{M_A}Y\}\ \text{mod,}\quad\text{mais} \& \{\underbrace{\complement_M X}_{X},\underbrace{\complement_{\unsure{M|m}}Y}_{Y}\}
\end{tikzcd}}
% L15–17
mod, donc comme \deleted{$A$ \unsure{s.}} les deux \unsure{termes} sont mod. avec $A$) $\{A\wedge X',A\wedge Y'\}$ modéré, OK.

% L18
5) Trivial

% L19–22
6) Soit $X\subset A$, les fermés dans $A$ sont les $Z=\deleted{\unsure{X|Z}}Y\wedge A$, avec $Y\in M_f$, $\{Y,A\}$ mod \deleted{\unsure{les} $\unsure{Z}$} \inserted{les $Z$} qui \unsure{contient|contiennent} $X$ \deleted{sont \unsure{ceux|les $Y$} qui \unsure{contiennent}} \inserted{\unsure{correspondent} \deleted{\unsure{à|i.e.} $Y\supset X$}} avec $Y$ t.q. $Y\supset X$ i.e. $Y\supset\overline X$, et le
